\documentclass[../../../include/open-logic-section]{subfiles}

\begin{document}

\olfileid{sth}{ord-arithmetic}{expo}
\olsection{Pemangkatan Ordinal}

Saiki kita ngalih menyang pemangkatan ordinal. Sayangé, ora ana
definisi sintetis kang \emph{gampang dibayangaké} kanggo
pemangkatan ordinal.

Pancen \emph{ana} definisi sintetis kang eksplisit. Iki salah
sijiné. Upamané $\text{finfun}(\alpha,\beta)$ iku himpunan kabèh
fungsi $f \colon \beta \to \alpha$ kang nyukupi yèn $\Setabs{\gamma \in
\beta}{f(\gamma) \neq 0}$ padha cacahé karo sawatara wilangan asli;
ing konstruksi iki dhasar ordinal kudu ora nol. Tetepna pengurutan
becik ing $\text{finfun}(\alpha,\beta)$ kanthi $f
\sqsubset g$ yen lan mung yen $f \neq g$ lan $f(\gamma_0) < g(\gamma_0)$, kanthi
$\gamma_0 = \text{max}\Setabs{\gamma \in \beta}{f(\gamma) \neq
g(\gamma)}$. Banjur kita bisa netepaké $\ordexpo{\alpha}{\beta}$
minangka $\ordtype{\text{finfun}(\alpha, \beta), \sqsubset}$. Potter
nganggo definisi eksplisit iki, banjur langsung njlentrehaké:
\begin{quote}
	Pamilihan pengurutan iki mung ditemtokaké déning karep kita
	kanggo olèh definisi pemangkatan ordinal kang manut syarat
	rekursif kang trep\ldots, lan iki luwih angel dibayangaké
	tinimbang gunggung terurut utawa ping-pingan terurut.
	\citep[kaca~199]{Potter2004}
\end{quote}
Pancen. Panambahan kita terangaké minangka “salinan $\alpha$
diterusaké salinan $\beta$”, lan ping-pingan minangka “urutan
$\beta$ salinan $\alpha$”. Nanging ora ana ukara ringkes kang
cocog kanggo nerangaké $\text{finfun}(\alpha, \gamma)$. Mula kita
bakal menehi definisi pemangkatan ordinal \emph{liwat} rekursi
transfinit, yaiku:

\begin{defn}\ollabel{ordexporecursion}
\begin{align*}
	\ordexpo{\alpha}{0} &= 1\\
	\ordexpo{\alpha}{\beta\ordplus 1} &=\ordexpo{\alpha}{\beta} \ordtimes \alpha\\
	\ordexpo{\alpha}{\beta} &= \bigcup_{\delta < \beta}\ordexpo{\alpha}{\delta}& & \text{yèn $\beta$ ordinal limit}
\end{align*}
\end{defn}

Yèn kita nyambut gawé \emph{minangka} ahli téori himpunan, kita
mbokmenawa kepéngin nliti sawetara sipat pemangkatan ordinal.
Nanging, ora akèh manèh kang perlu kita tambahi, kajaba nyathet
kasunyatan kang ora nggumunaké yèn pemangkatan ordinal ora
komutatif. Dadi $\ordexpo{2}{\omega} =
\bigcup_{\delta < \omega}\ordexpo{2}{\delta} = \omega$, déné
$\ordexpo{\omega}{2} = \omega \ordtimes \omega$. Nanging kita
ora kudu \emph{ngarep-arep} pemangkatan komutatif, amarga
pemangkatan wilangan asli uga ora komutatif: $\ordexpo{2}{3} = 8 < 9 = \ordexpo{3}{2}$.

\begin{prob}
Kanthi Induksi Transfinit, buktèkna yèn kanggo dhasar ordinal
kang ora nol, definisi
$\ordexpo{\alpha}{\beta} = \ordtype{\text{finfun}(\alpha, \beta),
\sqsubset}$ ngasilaké persamaan rekursi ing
\olref[sth][ord-arithmetic][expo]{ordexporecursion}.
\end{prob}

\end{document}
